\documentclass{book}
\usepackage[main=french,spanish,latin]{babel}
\usepackage[T1]{fontenc}
\usepackage{fontspec}
\usepackage{csquotes}
\usepackage[teiexport, divs=ekdosis, poetry=verse]{ekdosis}
\usepackage{setspace}
\usepackage{lettrine}
\usepackage{hyperref}
\usepackage{zref-user,zref-abspage}
% ... autres commandes utiles (optionnel)

\DeclareWitness{A}{1667}{Edition princeps}
\DeclareWitness{B}{1671}{Réédition de 1671}
\DeclareWitness{C}{1676}{Collective de 1676}
\DeclareWitness{D}{1687}{Collective de 1687}
\DeclareWitness{E}{1697}{Dernière du vivant de l'auteur}
\DeclareWitness{F}{2025}{Régularisée}
\begin{document}
\begin{ekdosis}
\begin{titlepage}
    \centering
    {\Huge \textbf{\MakeUppercase{Andromaque}} \par}
    \vspace{1.5cm}
    {\LARGE Jean Racine \par}
    \vspace{2cm}
    {\Large \textit{Édition critique par Clémentine Gheeraert} \par}
    \vspace{2cm}
    {\large \today}
    \vfill
\end{titlepage}
\ekddiv{type=scene, n=3, depth=3}
\stage{\app{
        \lem[wit={A, B, D, E, F}]{SCENE III.}
        \rdg[wit={C}]{SCENE. III.}
      }}
\stage{\app{
        \lem[wit={C, D, E, F}]{PYRRHUS, PHŒNIX.}
        \rdg[wit={A, B}]{PYRRHVS, PHŒNIX.}
      }}
    \begin{speech}
      \speaker{PHŒNIX.}
      \begin{ekdverse}
        \vnum{249}{
AInsi vous  
      \app{
        \lem[wit={C, D, E, F}]{l'envoyez}
        \rdg[wit={A}]{l’enuoyez}
        \rdg[wit={B}]{l'enuoyez}
      } aux pieds de sa  
      \app{
        \lem[wit={A, B, C, D, F}]{Maistresse~?}
        \rdg[wit={E}]{Maîtresse~?}
      }  \\    
        }
      \end{ekdverse}
    \end{speech}
    \begin{speech}
      \speaker{\app{
        \lem[wit={C, D, E, F}]{PYRRHUS.}
        \rdg[wit={A, B}]{PYRRHVS.}
      }}
      \begin{ekdverse}
        \vnum{250}{
On dit qu'il a long-temps brûlé pour la Princesse.  \\    
        }
      \end{ekdverse}
    \end{speech}
    \begin{speech}
      \speaker{PHŒNIX.}
      \begin{ekdverse}
        \vnum{251}{
Mais si ce feu, Seigneur, vient à se rallumer,  \\    
        }
        \vnum{252}{
S'il  
      \app{
        \lem[wit={A, B, C, E, F}]{luy}
        \rdg[wit={D}]{lui}
      } rendoit son  
      \app{
        \lem[wit={A, B, C, F}]{Cœur,}
        \rdg[wit={D, E}]{cœur,}
      } s'il s'en faisoit  
      \app{
        \lem[wit={A, B, F}]{aimer~?}
        \rdg[wit={C, D, E}]{aimer~!}
      }  \\    
        }
      \end{ekdverse}
    \end{speech}
    \begin{speech}
      \speaker{\app{
        \lem[wit={C, D, E, F}]{PYRRHUS.}
        \rdg[wit={A, B}]{PYRRHVS.}
      }}
      \begin{ekdverse}
        \vnum{253}{
Ah~! qu'ils s'aiment, Phœnix,  
      \app{
        \lem[wit={C, D, E, F}]{j'y}
        \rdg[wit={A, B}]{i'y}
      } consens. Qu'elle parte.  \\    
        }
        \vnum{254}{
Que  
      \app{
        \lem[wit={A, B, C, E, F}]{charmez}
        \rdg[wit={D}]{charmés}
      }  
      \app{
        \lem[wit={C, D, E, F}]{l'un}
        \rdg[wit={A, B}]{l'vn}
      } de l'autre, ils  
      \app{
        \lem[wit={B, C, D, E, F}]{retournent}
        \rdg[wit={A}]{retournẽt}
      } à Sparte.  \\    
        }
        \vnum{255}{
Tous nos Ports sont  
      \app{
        \lem[wit={C, D, E, F}]{ouverts}
        \rdg[wit={A, B}]{ouuerts}
      } \& pour elle \& pour  
      \app{
        \lem[wit={A, B, C, E, F}]{luy.}
        \rdg[wit={D}]{lui.}
      }  \\    
        }
        \vnum{256}{
Qu'elle m'épargneroit de contrainte \&  
      \app{
        \lem[wit={A, B, C, F}]{d'ennuy~!}
        \rdg[wit={D, E}]{d'ennui~!}
      }  \\    
        }
      \end{ekdverse}
    \end{speech}
    \begin{speech}
      \speaker{PHŒNIX.}
      \begin{ekdverse}
        \vnum{257.1}{
Seigneur....\\    
         }
      \end{ekdverse}
    \end{speech}
    \begin{speech}
      \speaker{\app{
        \lem[wit={C, D, E, F}]{PYRRHUS.}
        \rdg[wit={A, B}]{PYRRHVS.}
      }}
      \begin{ekdverse}
        \SetLineation{
                                lineation=none,
                                modulo,
                                vmodulo=0
                            }
        \vnum{257.2}{
\hspace*{5em}\app{
        \lem[wit={C, D, E, F}]{Une}
        \rdg[wit={A, B}]{Vne}
      } autre fois \app{
        \lem[wit={C, D, E, F}]{je}
        \rdg[wit={A, B}]{ie}
      } \app{
        \lem[wit={C, E, F}]{t'ouvriray}
        \rdg[wit={A, B}]{t'ouuriray}
        \rdg[wit={D}]{t'ouvrirai}
      } mon Ame,\\    
         }
        \resetvlinenumber[258]
        \SetLineation{vmodulo=5}
        \vnum{258.1}{
Andromaque \app{
        \lem[wit={A, B, C, F}]{paroist.}
        \rdg[wit={D}]{paroît.}
        \rdg[wit={E}]{paroist}
      }\\    
         }
      \end{ekdverse}
    \end{speech}

\end{ekdosis}
\end{document}
